\section{API 时代的 Agent：从通用函数到可执行系统}

\subsection{Universal Function 视角}
在 API-only 条件下，LLM 可以被视作“自然语言条件下的通用函数”。它在文本层面具备很强通用性，因此工程重心会转向“如何把这个函数嵌入可执行流程”。

\begin{importantbox}{API 不是叶子节点，而是控制器}
讲者强调：在现代 Agent 系统里，LLM API 不只是固定程序里的一个子调用；它常常直接决定执行流程、工具选择与中间状态更新，实际上在扮演控制器角色。
\end{importantbox}

\begin{figure}[H]
\centering
\includegraphics[width=0.93\textwidth]{frame-03-abstraction.jpg}
\caption{Agentic Abstraction：模型、工具、人类统一抽象}
\end{figure}
\noindent\footnotesize 画面证据：字幕约在 00:04:18--00:05:30，讲者讨论 “API controls execution flow” 与“memory stream / retrieval / action”。\normalsize

\subsection{记忆流、检索与反思循环}
讲者给出的基础 Agent 架构可以拆解为四个闭环部件：观察（observation）、记忆（memory stream）、检索（retrieval）、执行（action）。  
进一步地，系统会加入“反思”（reflection）机制，让模型对过去轨迹做高层总结，再把总结回灌到后续决策。

\begin{knowledgebox}{为什么 Reflection 在 Agent 中如此关键}
工具调用型 Agent 的失败往往不是单次推理错误，而是多步轨迹累积偏差。  
Reflection 的价值在于把“轨迹噪声”压缩成可复用的策略信号，使后续步骤不再盲目重复。
\end{knowledgebox}

\begin{warningbox}{记忆系统最容易出现的两个陷阱}
\begin{itemize}
    \item 无界累积：上下文持续膨胀导致检索质量下降；
    \item 错误固化：错误总结被重复引用，形成“自我强化幻觉”。
\end{itemize}
因此，记忆不是“存得越多越好”，而是“压缩质量与可验证性”更重要。
\end{warningbox}

\subsection{本章小结}
API 时代的 Agent 关键不是“多调用几个工具”，而是构建可控的执行闭环：让 LLM 成为控制器，同时用记忆与反思机制减少长链路误差。


